\section{Расширения сертификата открытого ключа}\label{FMT.Ext}

\subsection{Общие положения}\label{FMT.Ext.Intro} 

Расширения сертификата описывают дополнительную информацию о субъекте, эмитенте
или о самом сертификате. В зависимости от роли владельца сертификат может
содержать те или иные наборы расширений.

В соответствии с СТБ~34.101.19 расширение может быть обязательным или
необязательным, критическим или некритическим. Обязательное расширение~--- это
расширение, которое должно присутствовать в сертификате.  Критическое
расширение~--- это обязательное  расширение, которое должно быть корректным: при
нарушении корректности сторона, проверяющая сертификат, должна завершить
проверку с ошибкой.

Далее определяются допустимые расширения сертификатов, издаваемых КУЦ,
ИУЦ и ПУЦ. Если не оговорено противное, каждое из определяемых
расширений является обязательным для каждой роли владельца сертификата.

Сертификаты, издаваемые ПУЦ, могут содержать расширения, дополнительные к
перечисляемым ниже. Включение дополнительных расширений в сертификаты,
издаваемые КУЦ и ИУЦ, запрещено.

Названия расширений даются в соответствии с СТБ~34.101.19 (с прописной буквы).
Исключение составляют расширения~\code{ExtKeyUsage}
и~\code{AuthorityInfoAccess}, названия которых в СТБ~34.101.19 сопровождаются
суффиксом~\code{Syntax}.

\subsection{Расширения \code{SubjectKeyIdentifier} и
\code{AuthorityKeyIdentifier}}\label{FMT.Ext.SKID}

Расширения \code{SubjectKeyIdentifier} и \code{AuthorityKeyIdentifier} описывают
хэш-значения открытых ключей субъекта и эмитента сертификата соответственно.

Расширения являются обязательными, за исключением: \code{AuthorityKeyIdentifier}
может не включаться в самоподписанные сертификаты КУЦ. Расширения являются 
некритическими.

Хэш-значения должны вычисляться либо с помощью алгоритма~\code{belt-hash},
определенного в СТБ~34.101.31, либо с помощью алгоритма~SHA-1, определенного
в~\cite{SHA1}. В первом случае хэш-значение представляет собой строку из~$32$
октетов, во втором~--- строку из~$20$ октетов.

\begin{note}
Расширение \code{AuthorityKeyIdentifier} рекомендуется не включать в 
самоподписанные сертификаты КУЦ, сокращая тем самым размер сертификатов и 
упрощая их детектирование и проверку.
\end{note}

\begin{note}
Расширения~\code{SubjectKeyIdentifier} и \code{AuthorityKeyIdentifier} облегчают
построение цепочек сертификатов, не отвечая при этом за проверку цепочек.
Поэтому в расширениях разрешается использовать алгоритм хэширования SHA-1,
признанный на сегодняшний день криптографически нестойким. Разрешение на
использование SHA-1  продиктовано необходимостью обеспечивать совместимость с
действующими системами защиты информации. В тех случаях, когда совместимость не
нужна, следует использовать~\code{belt-hash}.
\end{note}

\subsection{Расширение \code{KeyUsage}}\label{FMT.Ext.KU}

Расширение~\code{KeyUsage} описывает назначение открытого ключа. Описание
представляет собой комбинацию флагов, определенных в СТБ~34.101.19.

Расширение является критическим.

В таблице~\ref{Table.FMT.Ext.KU} перечислены флаги \code{KeyUsage} для сторон
различных ролей.
%
Пропуск в таблице означает обязательное отсутствие флага, <<$+$>>~---
обязательное присутствие.

\begin{table}
\caption{Флаги \code{KeyUsage}}
\label{Table.FMT.Ext.KU}
\begin{tabular}{|l|c|c|c|c|c|}
\hline
\multirow{1}{*}[44pt]{Сторона} & 
\rotatebox{90}{\code{digitalSignature}~} &
\rotatebox{90}{\code{nonRepudiation}~} & 
\rotatebox{90}{\code{keyEncipherment}~} & 
\rotatebox{90}{\code{keyCertSign}~} & 
\rotatebox{90}{\code{cRLSign}~}\\
\hline
\hline
КУЦ          &   &   &   & + & + \\
\hline
ИУЦ  &       &   & + & + & + \\
\hline
ПУЦ          &   &   & + & + & + \\
\hline
ЦАC          & + &   &   &   & + \\
\hline
РЦ	         & + & + & + &   &   \\
\hline
OCSP-сервер  & + & + &   &   &   \\
\hline
СШВ          & + & + &   &   &   \\
\hline
СЗД          & + & + &   &   &   \\
\hline
СИ           & + & + & + &   &   \\
\hline
TLS-сервер   & + &   & + &   &   \\
\hline
ФЛ    	     & + & + & + &   &   \\
\hline                       
ЮП           & + & + & + &   &   \\
\hline                          
ЮЛ           & + & + & + &   &   \\
\hline
КА           & + &   & + &   &   \\
\hline                                     
\end{tabular}
\end{table}

\begin{note*}
Флаг \code{nonRepudiation} означает согласие с содержимым подписанного
сообщения, указывает на юридическую силу подписи. Чтобы подчеркнуть юридический 
аспект, в обновленных редакциях стандарта~\cite{X509} флаг \code{nonRepudiation} 
переименован в~\code{contentCommitment}.
\end{note*}

\subsection{Расширение \code{ExtKeyUsage}}\label{FMT.Ext.EKU}

Расширение \code{ExtKeyUsage} описывает область применения ключей 
сертификата. Описание представляет собой набор идентификаторов АСН.1. 

Расширение \code{ExtKeyUsage} не должно включаться в сертификаты УЦ, ЦАС и РЦ,
может включаться в сертификаты КА и должно включаться в сертификаты остальных
сторон. Расширение является критическим.

В сертификате OCSP-сервера расширение должно содержать
идентификатор \verb|id-kp-OCSPSigning|, определенный в СТБ~34.101.19.

В сертификате СЗД расширение должно содержать
идентификатор \verb|id-kp-dvcs|, определенный в СТБ~34.101.81.

В сертификате СШВ расширение должно содержать
идентификатор \verb|id-kp-timeStamping|, определенный в СТБ~34.101.19.

В сертификатах СИ и TLS-сервера расширение должно содержать
идентификатор \verb|id-kp-serverAuth|, определенный в СТБ~34.101.19.

В сертификатах ФЛ, ЮП и ЮЛ расширение должно содержать 
идентификаторы~\verb|id-kp-clientAuth| и~\verb|id-kp-emailProtection|,  
определенные в СТБ~34.101.19.

В сертификате стороны, которая выступает в роли сервера (клиента)
терминального режима, расширение \code{ExtKeyUsage} должно содержать
идентификатор~\code{bpki-eku-serverTM} (\code{bpki-eku-clientTM}). 
Идентификаторы определены в приложении~\ref{ASN1}.

\subsection{Расширение \code{CertificatePolicies}}\label{FMT.Ext.CP}

Расширение \code{CertificatePolicies} описывают политику, в соответствии 
с которой был выпущен сертификат, и цели, в которых сертификат может 
использоваться. 

Расширение~\code{CertificatePolicies} не должно включаться в сертификаты
КУЦ и должно включаться в сертификаты остальных сторон. 
Расширение является некритическим.

Описание политики состоит из пунктов, представленных идентификаторами 
АСН.1. В пунктах должны быть опущены опциональные классификаторы. 

В сертификатах ИУЦ и ПУЦ расширение должно содержать единственный 
пункт с идентификатором~\code{anyPolicy}, определенным в СТБ~34.101.19.

В сертификате конечного участника расширение 
должно содержать пункты с идентификаторами его ролей.
Идентификаторы определены в приложении~\ref{ASN1}
в соответствии с таблицей~\ref{Table.ENTITIES.Roles}. 
Расширение может включать пункты нескольких ролей.
Например, в сертификате оператора РЦ указываются две роли~--- ЮП и РЦ.

В сертификате конечного участника расширение~\code{CertificatePolicies} 
может содержать дополнительные пункты, отличные от пунктов 
ролей. Например, пункт политики, в соответствии с которой 
проверялось владение TLS-сервером заявленным DNS-именем.

\subsection{Расширение \code{BasicConstraints}}\label{FMT.Ext.BC}

Расширение~\code{BasicConstraints} дифференцирует сертификаты УЦ и
конечных участников. Дополнительно расширение ограничивает 
длину цепочек сертификатов, подчиненных сертификату УЦ.

Расширение \code{BasicConstraints} является критическим.

В сертификатах УЦ флаг~\code{cA} расширения должен быть 
установлен, в сертификатах конечных участников~--- сброшен. 

В сертификатах КУЦ, ИУЦ и конечных участников компонент 
\code{pathLenConstraint} должен отсутствовать,
а в сертификате ПУЦ~--- принимать значение~0. 
Это значение означает запрет на выпуск ПУЦ сертификатов 
другим УЦ. 

\subsection{Расширение \code{SubjectAltName}}\label{FMT.Ext.SAN}

Расширение~\code{SubjectAltName} содержит дополнительные 
идентификационные данные, описанные в~\ref{ENTITIES.SAN}. 

Расширение является некритическим и необязательным.

\subsection{Расширение \code{CRLDistributionPoints}}\label{FMT.Ext.CDP}

Расширение \code{CRLDistributionPoints} описывает расположение СОС, 
выпускаемых эмитентом сертификата. 

Расширение~\code{CRLDistributionPoints} 
не должно включаться в сертификаты КУЦ и должно включаться в
сертификаты остальных сторон. Расширение является некритическим и обязательным.

Отдельные точки распространения списков отзыва описываются
типом~\code{DistributionPoint}. Опциональные компоненты~\code{reasons} и
\code{cRLIssuer} этого типа должны быть опущены, а в
компоненте~\code{distributionPoint} должен быть указан URI-адрес точки
распространения (через выбор сначала варианта~\code{fullName}, а затем
варианта~\code{uniformResourceIdentifier}).

% info: адрес должен быть задан по схеме "http", схема "https" может 
% привести к циклической ссылке:
%   https://learn.microsoft.com/en-us/answers/questions/445425/

\subsection{Расширение~\code{AuthorityInfoAccess}}\label{FMT.Ext.AIA}

Расширение~\code{AuthorityInfoAccess} описывает информационные ресурсы, 
связанные с эмитентом сертификата.
%
Расширение состоит из пунктов типа~\code{AccessDescription}. 
Каждый такой пункт в свою очередь состоит из компонентов 
\code{accessMethod} (тип ресурса) и \code{accessLocation} 
(расположение ресурса). 

Расширение не должно включаться в сертификаты КУЦ и ИУЦ и должно 
включаться в сертификаты остальных сторон. Расширение является некритическим.

Первый пункт расширения описывает расположение сертификатов (одного или
нескольких) эмитента. В компоненте \code{accessMethod} пункта должен быть
установлен идентификатор~\verb|id-ad-caIssuers|, определенный в СТБ~34.101.19,
%
а в компоненте~\code{accessLocation}~--- URI-адрес сертификатов эмитента
(через выбор варианта~\code{uniformResourceIdentifier}).

Второй пункт расширения описывает расположение OCSP-сервера, к которому следует 
обратиться, чтобы получить информацию о статусе текущего сертификата.
В компоненте \code{accessMethod} пункта должен быть установлен 
идентификатор~\verb|id-ad-ocsp|, определенный в СТБ~34.101.19,
%
а в компоненте~\code{accessLocation}~--- URI-адрес OCSP-сервера
(через выбор варианта~\code{uniformResourceIdentifier}).
